% ReportKit minimal report skeleton (v1.1.0)
% Copy this file to start a new report. Sections are placeholders — remove
% any that don't fit the task per SKILL.md §3.
\documentclass{reportkit}

\setreportkitleftheader{REPORT FAMILY}
\setreportkitfooter{Short report descriptor}
\setreportkitversion{v1.0}

\title{Report title}
\author{Author or team}
\date{\today}

\begin{document}
\maketitle

\begin{execsummary}
State the problem, the strongest evidence, the conclusion, and the implication in a short narrative.
Avoid turning the executive summary into a list of disconnected highlights.
\end{execsummary}

\section{Problem}
Define the question, scope, and why it matters.

\begin{researchproblem}{Primary question}
What evidence would distinguish the competing explanations or options?
\end{researchproblem}

\section{Evidence and method}
Explain the evidence, data, method, and assumptions needed to interpret the result.

\begin{assumption}{Key assumption}
State a premise that materially affects the analysis.
\end{assumption}

\section{Results}
Present the minimum set of equations, tables, metrics, and figures needed to support the conclusion.

\begin{metric}{25,000}{commits per quarter}
Explain what this quantity measures and why it changes the reader's interpretation, decision, or
control requirement.
\end{metric}

\section{Implications}
Separate what the evidence shows from what should be done about it.

\begin{decisionpoint}{Decision required}
State the trade-off, options, criterion, and recommended default if appropriate.
\end{decisionpoint}

\section{Risks and limitations}
Describe meaningful uncertainty and failure modes rather than boilerplate caveats.

\begin{limitationnote}{Boundary condition}
State where the evidence should not be generalized.
\end{limitationnote}

\end{document}
